% Appendices for Dissertation
% These appendices support Chapter 3 (Materials and Methods)

\appendix

\chapter{Supplementary Materials}\label{app:supplementary}

%============================================================================
\section{Conversion and Judging Prompts}\label{app:conversion_prompts}

This section provides the complete prompts used in Phase 2 (MCQ to OSQ conversion) and Phase 5 (LLM-as-a-judge scoring).

\subsection{Phase 2 Stage 1: MCQ Suitability Classification Prompt}

\begin{lstlisting}[basicstyle=\small\ttfamily,breaklines=true]
You are an expert in educational assessment design. Your task is to evaluate whether a multiple-choice question (MCQ) is suitable for conversion to an open-ended short-answer question (OSQ).

Evaluate the following MCQ across five dimensions, scoring each from 1 (very poor) to 10 (excellent):

1. **Assessability**: Can the answer be objectively evaluated in open-ended format without relying on multiple-choice options?
2. **Canonical Answer Clarity**: Is there a well-defined, unambiguous correct answer?
3. **Ambiguity**: Would the question remain clear and unambiguous without the MCQ structure?
4. **Scope Appropriateness**: Is the expected answer reasonable in length for short-answer format (not requiring multi-paragraph essays)?
5. **Intent Preservation**: Can the original knowledge construct being tested be maintained without showing answer options?

MCQ to evaluate:
Question: {{question}}
A. {{choiceA}}
B. {{choiceB}}
C. {{choiceC}}
D. {{choiceD}}
Correct Answer: {{answer}}

Provide your evaluation in JSON format:
{
  "suitability_score": <1-10 overall score>,
  "justification": "<brief reasoning for overall score>",
  "dimension_scores": {
    "assessability": <1-10>,
    "canonical_answer": <1-10>,
    "ambiguity": <1-10>,
    "scope": <1-10>,
    "intent_preservation": <1-10>
  },
  "concerns": "<any specific concerns about conversion>"
}
\end{lstlisting}

\subsection{Phase 2 Stage 2: MCQ to OSQ Conversion Prompt}

\begin{lstlisting}[basicstyle=\small\ttfamily,breaklines=true]
You are an expert in educational assessment design. Your task is to convert a multiple-choice question (MCQ) to an open-ended short-answer question (OSQ) while preserving the original knowledge construct, difficulty level, and cognitive demand.

Guidelines:
1. Remove MCQ-specific artifacts (e.g., "which of the following", explicit A/B/C/D options)
2. Maintain the original intent (e.g., if testing recall, don't change to application)
3. Preserve any numerical calculations or technical specifications
4. Do NOT "improve" or make the question more challenging
5. Ensure the question can stand alone without needing the distractor options for context

MCQ to convert:
Question: {{question}}
A. {{choiceA}}
B. {{choiceB}}
C. {{choiceC}}
D. {{choiceD}}
Correct Answer: {{answer}} - {{correct_choice_text}}
Justification: {{justification}}

Create:
1. An open-ended question prompt (without showing answer choices)
2. A canonical expected answer based on the original correct choice
3. A detailed grading rubric with three levels:
   - Full credit: Criteria for complete correct answer
   - Partial credit: Criteria for partially correct answer
   - No credit: Conditions warranting no credit
4. Bloom's taxonomy classification for the cognitive level

Provide your conversion in JSON format:
{
  "osq_prompt": "<open-ended question>",
  "expected_answer": "<canonical correct answer>",
  "rubric": {
    "full_credit": "<specific criteria>",
    "partial_credit": "<specific criteria>",
    "no_credit": "<specific criteria>"
  },
  "blooms_level": "<Remember|Understand|Apply|Analyze|Evaluate|Create>",
  "blooms_justification": "<reasoning for taxonomy level>"
}
\end{lstlisting}

\subsection{Phase 5: Binary Judging Prompt}

\begin{lstlisting}[basicstyle=\small\ttfamily,breaklines=true]
You are grading a systems engineering question response.

Question: {{osq_prompt}}
Expected Answer: {{expected_answer}}

Student Response: {{model_response}}

Is the student response correct? Respond in JSON format:
{
  "correct": <true|false>,
  "reasoning": "<brief justification>"
}
\end{lstlisting}

\subsection{Phase 5: Rubric-Based Judging Prompt}

\begin{lstlisting}[basicstyle=\small\ttfamily,breaklines=true]
You are grading a systems engineering question response using a detailed rubric.

Question: {{osq_prompt}}
Expected Answer: {{expected_answer}}

Grading Rubric:
- Full credit (1.0): {{rubric_full_credit}}
- Partial credit (0.5): {{rubric_partial_credit}}
- No credit (0.0): {{rubric_no_credit}}

Student Response: {{model_response}}

Assign a score (0.0, 0.5, or 1.0) based on the rubric. Respond in JSON format:
{
  "score": <0.0|0.5|1.0>,
  "reasoning": "<explanation of which rubric criteria were met>",
  "rubric_level": "<full_credit|partial_credit|no_credit>"
}
\end{lstlisting}

\subsection{Phase 5: Multi-Dimensional Judging Prompt}

\begin{lstlisting}[basicstyle=\small\ttfamily,breaklines=true]
You are grading a systems engineering question response across five dimensions.

Question: {{osq_prompt}}
Expected Answer: {{expected_answer}}

Student Response: {{model_response}}

Evaluate the response on the following dimensions (0-10 scale each):

1. **Technical Accuracy** (0-10): Correctness of facts, principles, terminology, and technical details
2. **Conceptual Understanding** (0-10): Depth of systems thinking and grasp of underlying concepts
3. **Completeness** (0-10): Coverage of all required answer elements
4. **Clarity** (0-10): Communication quality, organization, and coherence
5. **Professional Relevance** (0-10): Real-world applicability and practical understanding

Respond in JSON format:
{
  "technical_accuracy": <0-10>,
  "conceptual_understanding": <0-10>,
  "completeness": <0-10>,
  "clarity": <0-10>,
  "professional_relevance": <0-10>,
  "aggregate_score": <0.0-1.0, computed as mean/10>,
  "reasoning": "<overall assessment explaining scores>"
}
\end{lstlisting}

\subsection{Phase 5: Chain-of-Thought Judging Prompt}

\begin{lstlisting}[basicstyle=\small\ttfamily,breaklines=true]
You are grading a systems engineering question response. Think step-by-step about the quality of the response.

Question: {{osq_prompt}}
Expected Answer: {{expected_answer}}

Student Response: {{model_response}}

Evaluate the response by reasoning through the following:
1. What are the key elements required in a correct answer?
2. Which of these elements are present in the student response?
3. Are there any errors, misconceptions, or missing critical details?
4. What are the specific strengths of the response?
5. What are the specific weaknesses of the response?
6. Overall, how well does this response answer the question?

Respond in JSON format:
{
  "reasoning_steps": [
    "<step 1 analysis>",
    "<step 2 analysis>",
    "<step 3 analysis>",
    ...
  ],
  "strengths": ["<strength 1>", "<strength 2>", ...],
  "weaknesses": ["<weakness 1>", "<weakness 2>", ...],
  "score": <0.0-1.0>,
  "justification": "<final summary explaining the score>"
}
\end{lstlisting}

%============================================================================
\section{Complete Model List}\label{app:model_list}

Table~\ref{tab:complete_model_list} lists all models evaluated in this research.

\begin{table}[htbp]
    \centering
    \caption{Complete list of evaluated models with key characteristics.}
    \label{tab:complete_model_list}
    \small
    \begin{tabular}{llrrlr}
        \toprule
        \textbf{Model Family} & \textbf{Model} & \textbf{Params (B)} & \textbf{VRAM (GB)} & \textbf{Type} & \textbf{Year} \\
        \midrule
        Llama 3.2 & llama3.2:3b & 3 & 6 & Open & 2024 \\
        Llama 3.2 & llama3.2:7b & 7 & 14 & Open & 2024 \\
        Llama 3.1 & llama3.1:70b & 70 & 43 & Open & 2024 \\
        Llama 3.1 & llama3.1:405b & 405 & 141 & Open & 2024 \\
        \addlinespace
        Qwen 2.5 & qwen2.5:7b & 7 & 14 & Open & 2024 \\
        Qwen 2.5 & qwen2.5:14b & 14 & 24 & Open & 2024 \\
        Qwen 2.5 & qwen2.5:32b & 32 & 38 & Open & 2024 \\
        Qwen 2.5 & qwen2.5:72b & 72 & 90 & Open & 2024 \\
        \addlinespace
        Mistral & mistral:7b-instruct & 7 & 14 & Open & 2023 \\
        Mistral & mixtral:8x7b & 47 & 48 & Open (MoE) & 2023 \\
        Mistral & mixtral:8x22b & 141 & 141 & Open (MoE) & 2024 \\
        \addlinespace
        DeepSeek & deepseek-r1:7b & 7 & 24 & Open (Reasoning) & 2024 \\
        DeepSeek & deepseek-r1:14b & 14 & 36 & Open (Reasoning) & 2024 \\
        \addlinespace
        QwQ & qwq:32b-preview & 32 & 48 & Open (Reasoning) & 2024 \\
        \addlinespace
        OpenAI & gpt-4o & N/A & Cloud & Proprietary & 2024 \\
        OpenAI & gpt-4o-mini & N/A & Cloud & Proprietary & 2024 \\
        OpenAI & gpt-3.5-turbo & N/A & Cloud & Proprietary & 2023 \\
        \addlinespace
        Anthropic & claude-3.5-sonnet & N/A & Cloud & Proprietary & 2024 \\
        Anthropic & claude-3-opus & N/A & Cloud & Proprietary & 2024 \\
        \addlinespace
        Google & gemini-1.5-pro & N/A & Cloud & Proprietary & 2024 \\
        Google & gemini-1.5-flash & N/A & Cloud & Proprietary & 2024 \\
        \bottomrule
    \end{tabular}
\end{table}

\textbf{Notes}:
\begin{itemize}
    \item \textbf{Params}: Billion parameters (approximate for MoE models)
    \item \textbf{VRAM}: Approximate GPU memory required for inference with 4-bit quantization
    \item \textbf{Type}: Open-source, proprietary, or specialized (MoE = Mixture of Experts, Reasoning = extended reasoning capabilities)
    \item \textbf{Year}: Release year
\end{itemize}

%============================================================================
\section{Dataset Visualizations}\label{app:dataset_viz}

This section presents additional visualizations of the SysEngBench dataset distribution across INCOSE knowledge areas.

\subsection{Sub-Category Distribution}

\begin{figure}[H]
    \centering
    % TODO: Include actual bar chart from Phase 1.2
    \includegraphics[width=\textwidth]{figures/subcategory_distribution.pdf}
    \caption{Horizontal bar chart showing distribution of questions across all 20 INCOSE sub-categories, sorted by question count.}
    \label{fig:subcategory_distribution}
\end{figure}

\subsection{Category Distribution}

\begin{figure}[H]
    \centering
    % TODO: Include actual pie chart from Phase 1.2
    \includegraphics[width=0.7\textwidth]{figures/category_pie.pdf}
    \caption{Pie chart showing percentage distribution across 6 major INCOSE categories.}
    \label{fig:category_pie}
\end{figure}

\subsection{Stacked Category Visualization}

\begin{figure}[H]
    \centering
    % TODO: Include actual stacked bar chart from Phase 1.2
    \includegraphics[width=\textwidth]{figures/category_stacked.pdf}
    \caption{Stacked bar chart showing hierarchical distribution of questions across categories and sub-categories.}
    \label{fig:category_stacked}
\end{figure}

%============================================================================
\section{Statistical Test Results Tables}\label{app:stat_tables}

This section provides detailed statistical test results referenced in Chapter~\ref{ch:materials_methods}.

\subsection{Position Bias Statistical Tests (Example)}

Table~\ref{tab:position_bias_stats_example} shows example statistical test results for position bias analysis (actual results in Chapter~\ref{ch:results}).

\begin{table}[H]
    \centering
    \caption{Example position bias statistical test results for selected models.}
    \label{tab:position_bias_stats_example}
    \small
    \begin{tabular}{lrrrrrr}
        \toprule
        \textbf{Model} & \textbf{Acc$_A$} & \textbf{Acc$_B$} & \textbf{Acc$_C$} & \textbf{Acc$_D$} & \textbf{$\chi^2$} & \textbf{$p$-value} \\
        \midrule
        llama3.2:3b & 0.623 & 0.611 & 0.605 & 0.598 & 2.14 & 0.543 \\
        qwen2.5:7b & 0.734 & 0.698 & 0.701 & 0.689 & 8.92 & 0.030* \\
        deepseek-r1:7b & 0.801 & 0.778 & 0.773 & 0.791 & 3.45 & 0.327 \\
        gpt-4o & 0.912 & 0.905 & 0.908 & 0.910 & 0.67 & 0.880 \\
        \bottomrule
    \end{tabular}
    \begin{tablenotes}
        \small
        \item * $p < 0.05$, ** $p < 0.01$, *** $p < 0.001$
        \item Chi-Square test for independence (df=3)
    \end{tablenotes}
\end{table}

\subsection{Format Comparison Statistical Tests (Example)}

Table~\ref{tab:format_comparison_stats_example} shows example statistical test results for MCQ vs. OSQ comparison (actual results in Chapter~\ref{ch:results}).

\begin{table}[H]
    \centering
    \caption{Example MCQ vs. OSQ statistical comparison for selected models.}
    \label{tab:format_comparison_stats_example}
    \small
    \begin{tabular}{lrrrrl}
        \toprule
        \textbf{Model} & \textbf{Acc$_{\text{MCQ}}$} & \textbf{Acc$_{\text{OSQ}}$} & \textbf{Diff} & \textbf{Cohen's $d$} & \textbf{Sig.} \\
        \midrule
        llama3.2:3b & 0.609 & 0.542 & +0.067 & 0.34 & * \\
        qwen2.5:7b & 0.706 & 0.673 & +0.033 & 0.18 & n.s. \\
        deepseek-r1:7b & 0.786 & 0.801 & -0.015 & -0.09 & n.s. \\
        gpt-4o & 0.909 & 0.894 & +0.015 & 0.12 & n.s. \\
        \bottomrule
    \end{tabular}
    \begin{tablenotes}
        \small
        \item Diff = Acc$_{\text{MCQ}}$ - Acc$_{\text{OSQ}}$
        \item Cohen's $d$: Small (0.2-0.5), Medium (0.5-0.8), Large ($\geq 0.8$)
        \item Sig.: * $p < 0.05$ (paired t-test), n.s. = not significant
    \end{tablenotes}
\end{table}

%============================================================================
\section{Example Conversions and Judging}\label{app:examples}

This section provides additional examples of MCQ to OSQ conversions and LLM-as-a-judge scoring.

\subsection{Example Conversion 1: Definitional Question}

\begin{tcolorbox}[title=Original MCQ]
\textbf{Question}: What is the primary purpose of a functional architecture in systems engineering?

\begin{enumerate}[label=\Alph*.]
    \item To document hardware components
    \item To define system functions and their relationships
    \item To specify manufacturing processes
    \item To establish project timelines
\end{enumerate}

\textbf{Correct Answer}: B
\end{tcolorbox}

\begin{tcolorbox}[title=Converted OSQ]
\textbf{OSQ Prompt}: What is the primary purpose of a functional architecture in systems engineering?

\textbf{Expected Answer}: To define system functions and their relationships

\textbf{Rubric}:
\begin{itemize}
    \item \textbf{Full credit}: Mentions defining/describing functions and their relationships/interactions/dependencies
    \item \textbf{Partial credit}: Mentions functions but not relationships, or vice versa
    \item \textbf{No credit}: Focuses on hardware, implementation, or unrelated concepts
\end{itemize}

\textbf{Bloom's Level}: Understand (comprehension of definition and purpose)
\end{tcolorbox}

\subsection{Example Judging: Rubric-Based Approach}

\begin{tcolorbox}[title=Question and Response]
\textbf{Question}: What is the primary purpose of a functional architecture in systems engineering?

\textbf{Expected}: To define system functions and their relationships

\textbf{Model Response}: "A functional architecture describes the functions that a system must perform and how those functions relate to each other and interact to achieve system objectives."
\end{tcolorbox}

\begin{tcolorbox}[title=LLM Judge Output (GPT-4o, Rubric-Based)]
\{\\
\quad "score": 1.0,\\
\quad "reasoning": "The response clearly identifies both key components required by the full credit rubric: (1) defining/describing system functions, and (2) their relationships/interactions. The response goes beyond minimum requirements by also mentioning system objectives, demonstrating comprehensive understanding.",\\
\quad "rubric\_level": "full\_credit"\\
\}
\end{tcolorbox}

%============================================================================
\section{Code and Data Availability}\label{app:code}

All code, data, and evaluation outputs supporting this dissertation are publicly available:

\subsection{GitHub Repository}

\textbf{URL}: \url{https://github.com/ryan-a-bell/dissertation}

\textbf{Contents}:
\begin{itemize}
    \item \texttt{src/phase1\_prep/}: Phase 1 data download and characterization notebooks
    \item \texttt{src/phase2\_conversion/}: Phase 2 MCQ to OSQ conversion pipeline
    \item \texttt{src/phase3\_variants/}: Phase 3 position variant generation
    \item \texttt{src/phase4\_inference/}: Phase 4 inference scripts and task configurations
    \item \texttt{src/phase5\_llm\_as\_a\_judge/}: Phase 5 judging scripts
    \item \texttt{src/phase6\_analysis/}: Phase 6 comprehensive analysis scripts
    \item \texttt{manuscript/}: LaTeX dissertation source files
\end{itemize}

\subsection{HuggingFace Datasets}

\begin{itemize}
    \item \textbf{SysEngBench (Base)}: \url{https://huggingface.co/datasets/rabell/SysEngBench}
    \item \textbf{SysEngBench-OSQ}: \url{https://huggingface.co/datasets/rabell/SysEngBench-OSQ}
    \item \textbf{Position Variants}: Included in base dataset repository
\end{itemize}

\subsection{Model Results Archive}

Due to large file sizes, complete evaluation outputs (JSON results, JSONL samples) are archived separately:

\textbf{Zenodo DOI}: [TO BE ASSIGNED UPON PUBLICATION]

%============================================================================
\section{Software Versions and Environment}\label{app:environment}

\subsection{Python Environment}

\begin{lstlisting}[language=bash,basicstyle=\small\ttfamily]
# Core packages
python==3.11.6
pandas==2.1.4
numpy==1.24.3
scipy==1.11.4
scikit-learn==1.3.2
matplotlib==3.8.2
seaborn==0.13.0

# HuggingFace ecosystem
datasets==2.14.6
transformers==4.36.0

# Evaluation
lm-eval==0.4.0
ollama==0.3.3
openai==1.6.1

# Utilities
tqdm==4.66.1
python-dotenv==1.0.0
jupyter==1.0.0
jupyterlab==4.0.9
\end{lstlisting}

\subsection{Hardware Specifications}

\paragraph{RunPod GPU Instances}
\begin{itemize}
    \item GPU: NVIDIA A40 (48GB VRAM) or multi-GPU configurations
    \item CPU: AMD EPYC 7532 (32 cores)
    \item RAM: 128GB DDR4
    \item Storage: 500GB NVMe SSD
    \item Network: 10 Gbps
\end{itemize}

\paragraph{Development Workstation}
\begin{itemize}
    \item OS: Ubuntu 22.04 LTS
    \item CPU: Intel Core i9-12900K
    \item RAM: 64GB DDR5
    \item GPU: NVIDIA RTX 3090 (24GB VRAM, for development/testing)
\end{itemize}
